Large language models are increasingly deployed for creative tasks such as brainstorming, recommendation, and cultural curation, yet their capacity for genuinely novel judgment remains poorly understood. We introduce the ``most underrated X'' question as a diagnostic for LLM novelty limitations: a task that requires second-order reasoning about the gap between quality and recognition. We query five LLMs across four model families on 40 categories spanning eight domains, collecting 1,716 responses. We find that models overwhelmingly converge on the same ``obviously underrated'' answers---the same picks that dominate Reddit threads and online forums. The mean intra-model agreement rate is 68.2\% (Cohen's $d = 2.42$ vs.\ chance, $p < 0.0001$), and 49.4\% of LLM top answers match Reddit consensus picks. We formalize this as the \textit{novelty paradox}: when multiple independent LLMs name the same item as ``most underrated,'' that item is, by definition, the consensus choice rather than a genuinely novel one. GPT-4.1 is the most predictable model (75\% Reddit match), while Llama~4 Maverick shows the most independence (20\% Reddit match). These results establish that current LLMs conflate ``frequently discussed as underrated'' with ``genuinely underrated,'' revealing a structural limitation in their capacity for original evaluative judgment.
